\section{Evaluation Plan and Established Evidence}
\label{sec:evaluation}

The evaluation distinguishes semantic proof, compilation, execution validation,
and cost. A successful native test is not a universal correctness proof, and
proof counts are not a measure of practical optimization value.

\subsection{Current Artifact Evidence}

For the loaded-root-with-offset compiler, the existing report records six
configurations and 708 full assembly calls. Separate path probes record 236
Clight calls and 21 probes of unmodified assembly. These cover candidate
acceptance and fallback, changing bounds through aliases, empty paths,
multi-array bodies, and repeated rewrites. The reports are bound to their
compiler/source digests; these are prior execution results, not a new matrix run
for this manuscript.

The earlier Figure~2 coverage report records sixteen unchanged-source calls;
it remains a fixed record, not an optimization-success count. The new frontend
closure covers 20 endpoints: 17 language results, one domain result, and two
fixtures, with 569 required dependencies and 1,104 source digests. It binds the
parent nested compiler closure, whose 25 endpoints connect original-source
checking, canonical execution, alias/candidate validation, kernel preservation,
and installation. The kernel is unchanged; the new compiler retains the same
42-global assumption set, with no additional global axiom.

The new extracted frontend compiler was run on the same adapted C source.
Six configurations---disabled, identity, interchange, $2\times3$ tiling, wrong
reindexing, and invalid domain---give 48 complete assembly calls. A further 27
calls use a harness with one additional input and unchanged kernel/context
prefix. All 75 calls compare all 2,048 output words, both headers, and public
counter exits against a GCC wrapping-arithmetic reference and an independent
source-word model. Separate instrumented Clight runs give sixteen dispatch
observations: two fast and six fallback calls for each reordered configuration.
These Clight observations are not assembly-path evidence.

Six debugger probes observe unmodified assembly. Three watch selected output
locations and distinguish their write order for one independent-array input
(Table~\ref{tab:nested-order}); the
other three check first-header or second-header-region alias fallback. The
second-header-region probe excludes a write that the cached iteration domain
would have made. Public exits and full outputs match source behavior. Probe
runs are separate from the 75-call assembly matrix. Exact commands, source and
compiler bindings, report hashes, and historical stages are in the evidence map.

\begin{table}[t]
\centering
\caption{Write order at three watched output locations in assembly for one
independent-array input. All three runs finish with public counters (3, 4, 5).}
\label{tab:nested-order}
\begin{tabular}{llr}
\toprule
Configuration & Output word indices & Function bytes \\
\midrule
Disabled & 15, 80, 160 & 161 \\
Interchange & 80, 160, 15 & 564 \\
2-by-3 tiling & 80, 15, 160 & 693 \\
\bottomrule
\end{tabular}
\end{table}
The identity configuration uses 528 function bytes. Code size is an observation
of the current generated programs, not a timing or profitability result.

\subsection{Experiments Needed for the Final Paper}

\paragraph{Semantic coverage.}
Classify accepted candidates, runtime refusals, and static matcher/validator
refusals. Use the same source for independent arrays, header/data aliasing,
empty outer and inner loops, overflowing offset expressions, limited readable
allocations, invalid candidate domains, and conflicting dependencies. Verify
complete arrays and public loop exits, and inspect executed store order to
distinguish a real fast path from unchanged-source compilation.

\paragraph{Condition construction.}
Compare the scan and compact sufficient conditions on a common source class.
Report accepted input domains and refused cases. Count captures, pointer
comparisons, numeric operations, and iterations separately. Existing work
counts show that loop-shaped guard code can still perform expensive repeated
comparisons; code-size reduction does not establish cheaper checking.

\paragraph{Runtime and compilation cost.}
Use the same CompCert version, target, and backend configuration for the original
and transformed programs. Measure check-only work where separable and total
execution cost, including refusal and empty paths. Report machine code size,
compile time, measurement environment, raw batches, and dispersion. Avoid
concurrent proof/build work during timing. Slowdowns and narrow acceptance are
results to explain rather than cases to remove.

\paragraph{Proof reuse.}
For multiple transformations, identify the shared kernel, language lemmas, and
domain services. List the site-specific assumptions and source/model bridges
still supplied by the author. Compare a common example against the nearest
existing interfaces before claiming reduced proof burden.

\pending{No timing, speedup, final acceptance fraction, or comparative author
effort has been measured for this manuscript. Extend this frontend's coverage to
multi-array read/write bodies, dependency refusals, same-allocation slices,
nonzero entry counters, and repeated/control contexts; validate a fresh-build
bootstrap. Measure compact conditions and author obligations on common examples.}
